\documentclass[11pt,a4paper]{article}

\usepackage[margin=1in]{geometry}
\usepackage{amsmath,amssymb,amsfonts}
\usepackage{mathtools}
\usepackage{lmodern}
\usepackage[T1]{fontenc}
\usepackage[utf8]{inputenc}
\usepackage{textcomp}
\usepackage{microtype}
\usepackage{parskip}
\usepackage{enumitem}
\usepackage{array}
\usepackage{booktabs}
\usepackage{hyperref}
\hypersetup{hidelinks,
  pdftitle={From Primitives to Bloch's Theorem},
  pdfauthor={Allen Proxmire}}

\setlength{\parindent}{0pt}
\setlength{\parskip}{6pt plus 2pt minus 1pt}

\title{\vspace{-1cm}\Large\bfseries From Primitives to Bloch's Theorem\\[6pt]
\large\mdseries A Walkthrough of the Event Density Derivation}
\author{Allen Proxmire}
\date{May 2026}

\begin{document}
\maketitle
\thispagestyle{empty}

\section{The Question}

\subsection*{What this walkthrough derives}

This walkthrough derives, from substrate primitives, the structural backbone of band theory:

\begin{enumerate}[noitemsep]
  \item \textbf{Bloch's theorem}: in a periodic rule-type substrate, the eigenchannels of the chain's evolution operator have the form
  \[
  \psi_k(x) = e^{ikx} u_k(x), \qquad u_k(x+a) = u_k(x),
  \]
  where $a$ is the spatial period of the substrate and $k$ is a quasi-momentum defined modulo $2\pi/a$.
  \item \textbf{The Brillouin zone}: the fundamental domain of $k$ values, $k \in (-\pi/a, \pi/a]$, forced by the discrete translation symmetry of the periodic substrate.
  \item \textbf{The band eigenvalue problem}: the periodic factor $u_{n,k}(x)$ satisfies a parameter-dependent eigenvalue problem on the unit cell, indexed by band number $n$ and quasi-momentum $k$, producing dispersion relations $E_n(k)$.
  \item \textbf{Band gaps}: forced by rule-type incompatibility within the unit cell --- frequency ranges where no Bloch eigenchannel exists.
  \item \textbf{Berry connection and curvature on the Brillouin zone}: $\mathcal{A}_n(k) = i\langle u_{n,k}|\partial_k u_{n,k}\rangle$ and $\Omega_n(k) = \partial_k \mathcal{A}_n(k)$, giving the geometric structure of the band manifold.
\end{enumerate}

Each is derived, not posited. The walkthrough is fully self-contained.

\subsection*{Why standard QM/solid-state physics treats Bloch's theorem as a formal eigenvalue result}

In the standard textbook derivation, one writes a single-particle Hamiltonian $H = -\hbar^2\nabla^2/(2m) + V(x)$ with $V(x+a) = V(x)$, observes that $H$ commutes with the translation operator $T_a$, and applies the spectral theorem to find simultaneous eigenfunctions. Bloch's theorem follows as a consequence of representation theory of the abelian group $\mathbb{Z}$ (translations by integer multiples of $a$). Band structure follows from solving the resulting eigenvalue problem on the unit cell.

This is mathematically correct. It is also mechanistically opaque. Standard treatments do not say \emph{what gives the periodic potential its substrate-level meaning}, \emph{why the chain must respect the substrate's translation symmetry}, \emph{what physical object the quasi-momentum $k$ describes at a deeper ontological level}, or \emph{what makes a band gap structurally forbidden rather than merely energetically suppressed}. The Hamiltonian's periodicity is given as a postulate; the theorem is its formal consequence.

\subsection*{What Event Density claims}

Bloch's theorem is FORCED at the substrate level of the Event Density (ED) framework whenever a chain propagates through a periodic rule-type substrate. The mechanism is structural: a periodic substrate generates a translation symmetry on the rule-type structure itself, the chain's evolution operator must commute with the substrate's translation operator (because both act on the same substrate-level rule structure), and the simultaneous eigenchannels of these commuting operators have the Bloch form by the algebraic structure of the abelian translation group. Band gaps emerge as ranges of energy for which no rule-type-compatible eigenchannel exists across the periodic unit cell --- substrate-level rule-type incompatibility, not dynamical suppression.

The substrate-level mechanism, in one sentence: \textbf{a periodic rule-type substrate forces the chain to align its identity with the periodicity, producing Bloch eigenchannels and a band structure indexed by quasi-momentum}.

\subsection*{The chain in summary}

The derivation chain runs:

substrate primitives with spatial periodicity (\S2) $\to$ periodic rule + translation operator + commutation with evolution operator (\S3) $\to$ Bloch form forced by commuting-operator eigenbasis (\S4) $\to$ band eigenvalue problem on unit cell (\S5) $\to$ reciprocal lattice + Brillouin zone geometry (\S6) $\to$ Berry connection and curvature on the band manifold (\S7) $\to$ substrate-level reading (\S8) $\to$ forced/inherited/open accounting (\S9) $\to$ exact claims (\S10).

The substrate primitives are five: participation rule with spatial periodicity, identity alignment under spatial translation, channel propagation through periodic gradients, translation operator as substrate-level rule-shift, and parameter-space bundle over the Brillouin zone.

\section{The Primitives}

\subsection{P-BL-1. Participation rule with spatial periodicity}

A \emph{participation rule} $r(x)$ is the substrate-level identity-encoding of a chain at spatial position $x$ --- the rule that determines how the chain interacts with ED-gradients at $x$. A \emph{periodic} participation rule satisfies
\[
r(x + a) = r(x) \quad \text{for all } x \in \mathbb{R},
\]
where $a > 0$ is the spatial period. The rule is single-typed at each $x$: at any one substrate-tick, the chain at position $x$ commits to exactly one rule from the alignment set $\mathcal{R}(x)$, and the alignment-set structure $\mathcal{R}(x)$ itself is periodic: $\mathcal{R}(x + a) = \mathcal{R}(x)$.

\textbf{Discrete translation symmetry.} The substrate-level statement: a translation by any integer multiple of $a$ leaves the rule structure invariant. Equivalently, two positions $x$ and $x + na$ (for $n \in \mathbb{Z}$) are substrate-levelly indistinguishable in their rule-type content.

\textbf{Algebraic structure.} As in the non-periodic case, the alignment set carries an inner-product structure $\langle r_i(x)|r_j(x)\rangle$ with $\langle r_i(x)|r_i(x)\rangle = 1$. Periodicity means $\langle r_i(x+a)|r_j(x+a)\rangle = \langle r_i(x)|r_j(x)\rangle$ for all $x$.

\textbf{Smooth dependence.} The rule depends smoothly on $x$ within each period, so $\nabla_x|r_n(x)\rangle$ is well-defined.

\subsection{P-BL-2. Identity alignment under spatial translation}

The \emph{identity alignment} of a chain at position $x$ is its commitment to one specific rule $r_n(x)$ within the alignment set $\mathcal{R}(x)$. Under a spatial translation $x \mapsto x + a$, the alignment-set structure is invariant (P-BL-1), so the chain's identity-commitment translates as
\[
r_n(x) \mapsto r_n(x + a) = r_n(x).
\]
The chain at $x + a$ aligns with the \emph{same labeled rule} as the chain at $x$, modulo the smooth $x$-dependence within each unit cell. This is the substrate-level statement of translation-equivariant identity preservation.

\subsection{P-BL-3. Channel propagation through periodic gradients}

A \emph{channel} is a participation pathway through ED-gradients carrying alignment threads. In a periodic rule-type substrate, the channel propagates through ED-gradients whose structure is periodic with period $a$. The channel's pre-individuation amplitude $\psi(x)$ is a complex-valued function on $\mathbb{R}$ (the substrate-level pre-image of the wavefunction).

\textbf{Evolution operator.} The chain's evolution is governed by a substrate-level evolution operator $H$ whose coarse-graining produces a Hamiltonian on the channel's Hilbert space. The eigenvalue equation
\[
H \psi(x) = E \psi(x)
\]
determines the channel's eigen-energies and eigenchannels.

\textbf{Periodicity-respecting evolution.} Because the substrate's rule structure is periodic with period $a$, the evolution operator $H$ inherits this periodicity: $H$ acts identically at $x$ and at $x + a$. This is the substrate-level statement that makes $H$ commute with the translation operator (\S3.2).

\subsection{P-BL-4. Translation operator $T_a$ as substrate-level rule-shift}

Define the \emph{translation operator} $T_a$ by
\[
(T_a \psi)(x) \equiv \psi(x - a).
\]
(Equivalently $T_a$ shifts the channel forward by $a$; some texts use the opposite sign convention. We use the convention where $T_a$ acts on amplitudes by pulling back the argument by $a$.)

\textbf{Unitarity.} $T_a$ preserves inner products: $\int|T_a\psi|^2 dx = \int|\psi(x-a)|^2 dx = \int|\psi(y)|^2 dy$ under change of variable $y = x - a$. Therefore $T_a^\dagger T_a = I$.

\textbf{Group structure.} Translations form an abelian group: $T_a T_b = T_{a+b}$, $T_a T_a^{-1} = I$, $T_a T_b = T_b T_a$. The discrete subgroup generated by $T_a$ is $\{T_{na} : n \in \mathbb{Z}\} \cong \mathbb{Z}$.

\textbf{Substrate-level meaning.} $T_a$ is the substrate-level shift of the rule structure by one period. P-BL-1 (substrate periodicity) is the statement that $T_a$ leaves the rule structure invariant.

\subsection{P-BL-5. Parameter-space bundle over the Brillouin zone}

The \emph{Brillouin zone} $\mathcal{B} = (-\pi/a, \pi/a]$ is the fundamental domain of quasi-momentum values, with the endpoints $-\pi/a$ and $\pi/a$ identified to make $\mathcal{B}$ topologically a circle. Each point $k \in \mathcal{B}$ carries a fibre Hilbert space spanned by the periodic factors $\{u_{n,k}(x)\}_n$ (derived in \S5).

\textbf{Bundle structure.} The total space is the disjoint union $E = \bigsqcup_{k \in \mathcal{B}} \mathcal{H}(k)$, where $\mathcal{H}(k)$ is the Hilbert space of $a$-periodic functions on the unit cell (defined precisely in \S5.1). The projection $\pi: E \to \mathcal{B}$ sends each $u_{n,k}$ to its base point $k$.

\textbf{Generally non-trivial.} The choice of phase for each $u_{n,k}$ may not extend to a globally smooth choice across all of $\mathcal{B}$ (especially across the BZ boundary $k = \pm\pi/a$). This non-triviality is what makes the Berry phase across the Brillouin zone potentially non-zero (\S7).

\textbf{No new substrate primitives are introduced beyond P-BL-1 through P-BL-5.}

\section{Periodic Rule-Type Structure and the Translation Operator}

\subsection{Periodic Hamiltonian}

The chain's evolution operator $H$ acts on the pre-individuation amplitude $\psi(x)$ via a kinetic term plus a periodic rule-type coupling:
\[
H = -\frac{\hbar^{2}}{2m} \frac{\partial^{2}}{\partial x^{2}} + V(x), \qquad V(x + a) = V(x),
\]
where $V(x)$ is the coarse-grained image of the periodic rule-type structure (P-BL-1), and the coefficient $-\hbar^2/(2m)$ is inherited from the standard kinetic-energy form (re-derived as follows: the substrate-level free-channel Hamiltonian acting on $\psi(x)$ is the second-derivative operator with coefficient $-\hbar^2/(2m)$; this is the coarse-graining of the V1 kernel's leading-order finite-width contribution, with the standard normalization).

Note: the choice of free-channel kinetic term is non-essential to the Bloch derivation. The structural argument requires only that $H$ contain the periodic potential and that $H$ act locally enough to commute with $T_a$ --- the kinetic-energy form serves as a concrete instantiation but the theorem applies to any local periodic Hamiltonian.

\subsection{Commutation $[H, T_a] = 0$}

We show that $H$ commutes with $T_a$.

The kinetic term $-\hbar^2/(2m)\,\partial^2/\partial x^2$ is translation-invariant: translation in the argument commutes with differentiation in the argument. Explicitly, $(T_a\partial^2\psi)(x) = (\partial^2\psi)(x-a) = \partial^2[\psi(x-a)] = \partial^2(T_a\psi)(x)$, so $T_a \circ \partial^2 = \partial^2 \circ T_a$.

The potential term $V(x)$ acts by multiplication. Compute:
\begin{align*}
(T_a V \psi)(x) &= (V\psi)(x - a) = V(x - a) \psi(x - a) = V(x) \psi(x - a) \\
&= V(x) (T_a \psi)(x) = (V T_a \psi)(x),
\end{align*}
using $V(x-a) = V(x)$ (periodicity). Therefore $T_a V = V T_a$, and combining with $T_a \partial^2 = \partial^2 T_a$:
\[
[H, T_a] = 0.
\]
This is the substrate-level statement that the chain's evolution respects the substrate's discrete translation symmetry (P-BL-3 + P-BL-4).

\subsection{Simultaneous eigenbasis}

By the spectral theorem (re-derived briefly: for any two commuting normal operators on a Hilbert space, there exists a basis of simultaneous eigenvectors --- the proof uses the fact that the eigenspaces of one operator are invariant under the other, and the second operator can be diagonalized within each eigenspace), $H$ and $T_a$ admit a simultaneous eigenbasis. We seek functions $\psi(x)$ satisfying
\[
H \psi(x) = E \psi(x) \quad \text{and} \quad T_a \psi(x) = \lambda \psi(x)
\]
simultaneously.

The structural content of Bloch's theorem lives entirely in the second equation: the eigenvalues of $T_a$ determine the form of $\psi(x)$. We address this in \S4.

\section{Deriving the Bloch Form}

\subsection{Eigenvalues of the translation operator}

We solve $T_a \psi(x) = \lambda \psi(x)$, i.e., $\psi(x - a) = \lambda \psi(x)$.

Because $T_a$ is unitary (\S P-BL-4), its eigenvalues lie on the unit circle: $|\lambda| = 1$. Write
\[
\lambda = e^{-ika} \quad \text{for some real } k \in \mathbb{R}.
\]
(The minus sign in the exponent is a convention chosen so that $\psi(x - a) = e^{-ika}\psi(x)$ implies $\psi(x + a) = e^{+ika}\psi(x)$, matching the standard Bloch convention.)

\subsection{The Bloch ansatz}

Define
\[
u_k(x) \equiv e^{-ikx} \psi(x),
\]
so that $\psi(x) = e^{ikx} u_k(x)$. Compute the action of $T_a$ on this form:
\[
(T_a \psi)(x) = \psi(x - a) = e^{ik(x-a)} u_k(x - a) = e^{-ika} \cdot e^{ikx} u_k(x - a).
\]
Comparing with $(T_a\psi)(x) = e^{-ika}\psi(x) = e^{-ika} e^{ikx} u_k(x)$:
\[
e^{ikx} u_k(x - a) = e^{ikx} u_k(x),
\]
so
\[
u_k(x - a) = u_k(x),
\]
equivalently $u_k(x + a) = u_k(x)$. \textbf{The function $u_k$ is periodic with the substrate period $a$.}

\subsection{The Bloch theorem statement}

We have shown: any simultaneous eigenchannel of $H$ and $T_a$ has the form
\[
\psi_k(x) = e^{ikx} u_k(x), \qquad u_k(x + a) = u_k(x).
\]
This is Bloch's theorem, derived from the substrate-level commutation $[H, T_a] = 0$ + the unitarity of $T_a$ + the periodicity of the rule-type structure.

\subsection{Periodicity of $k$ modulo $2\pi/a$}

The quasi-momentum $k$ enters only through the phase factor $e^{-ika}$ (equivalently $e^{ika}$). Two values $k$ and $k + 2\pi/a$ give the same eigenvalue of $T_a$:
\[
e^{-i(k + 2\pi/a)a} = e^{-ika - 2\pi i} = e^{-ika}.
\]
Therefore $k$ is defined modulo $2\pi/a$. The fundamental domain is the \textbf{Brillouin zone}
\[
\mathcal{B} = (-\pi/a, \pi/a],
\]
with the endpoints identified ($-\pi/a \sim \pi/a$). Topologically, $\mathcal{B}$ is a circle of circumference $2\pi/a$.

\textbf{Substrate-level statement.} The Brillouin zone is the parameter-space domain forced by the discrete translation symmetry of the substrate. The compactness of $\mathcal{B}$ (as a circle) is the substrate-level statement that translations by integer multiples of $a$ are substrate-levelly identified --- there is no substrate-level distinction between $k$ and $k + 2\pi/a$.

\subsection{The full Bloch eigenchannel}

Combining: any simultaneous eigenchannel of $H$ and $T_a$ is of the form
\[
\psi_{n,k}(x) = e^{ikx} u_{n,k}(x), \qquad u_{n,k}(x + a) = u_{n,k}(x), \qquad k \in \mathcal{B},
\]
with the band index $n$ labeling distinct eigenchannels at fixed $k$ (introduced in \S5).

\section{The Band Eigenvalue Problem}

\subsection{Reduction to the unit cell}

Insert the Bloch form $\psi_k(x) = e^{ikx} u_k(x)$ into the eigenvalue equation $H\psi = E\psi$:
\[
\left[-\frac{\hbar^{2}}{2m} \frac{\partial^{2}}{\partial x^{2}} + V(x)\right] [e^{ikx} u_k(x)] = E\, e^{ikx} u_k(x).
\]
Compute the kinetic term:
\begin{align*}
\partial_x [e^{ikx} u_k] &= e^{ikx} [ik u_k + \partial_x u_k], \\
\partial^{2}_x [e^{ikx} u_k] &= e^{ikx} [(ik)^{2} u_k + 2(ik) \partial_x u_k + \partial^{2}_x u_k] \\
&= e^{ikx} [-k^{2} u_k + 2ik \partial_x u_k + \partial^{2}_x u_k].
\end{align*}
So
\[
-\frac{\hbar^{2}}{2m} \partial^{2}_x [e^{ikx} u_k] = e^{ikx} \cdot \frac{\hbar^{2}}{2m} \cdot \bigl(-(\partial_x + ik)^{2}\bigr) u_k.
\]
Setting this in the eigenvalue equation and dividing by $e^{ikx}$:
\[
\left[\frac{\hbar^{2}}{2m} \bigl(-(\partial_x + ik)^{2}\bigr) + V(x)\right] u_k(x) = E\, u_k(x),
\]
equivalently
\[
H_k u_k(x) = E_k u_k(x), \qquad H_k \equiv -\frac{\hbar^{2}}{2m} (\partial_x + ik)^{2} + V(x).
\]
\textbf{This is the band eigenvalue problem.} $H_k$ acts on $a$-periodic functions $u_k(x)$ on the unit cell $[0, a)$. It is parameter-dependent in $k$.

\subsection{The unit-cell Hilbert space}

Define $\mathcal{H}_{\mathrm{cell}}$ as the space of square-integrable $a$-periodic complex functions on $[0, a)$ with inner product
\[
\langle u | v\rangle_{\mathrm{cell}} = \frac{1}{a} \int_0^a u^*(x)\, v(x)\, dx.
\]
The operator $H_k$ acts on $\mathcal{H}_{\mathrm{cell}}$. It is Hermitian with respect to this inner product (the kinetic term $-(\partial_x + ik)^2$ is the square of a Hermitian operator on $a$-periodic functions, and $V(x)$ is real).

\subsection{Discrete spectrum and band structure}

For each $k \in \mathcal{B}$, $H_k$ is a Hermitian operator on $\mathcal{H}_{\mathrm{cell}}$. Because $\mathcal{H}_{\mathrm{cell}}$ is (when restricted to functions with bounded gradient) effectively finite-dimensional in any finite-energy window, the spectrum of $H_k$ is discrete:
\[
H_k u_{n,k}(x) = E_n(k) u_{n,k}(x), \qquad n = 1, 2, 3, \ldots,
\]
with $E_1(k) \leq E_2(k) \leq \ldots$ ordered by increasing eigen-energy. The functions $u_{n,k}(x)$ are orthonormal:
\[
\langle u_{m,k} | u_{n,k}\rangle_{\mathrm{cell}} = \delta_{mn}.
\]
The dispersion relation $E_n(k)$ for fixed $n$ as $k$ varies over $\mathcal{B}$ is the \emph{$n$-th band}. The set $\{E_n(k) : k \in \mathcal{B}\}$ is a closed interval (the $n$-th band's energy range) by continuity of $H_k$ in $k$ + compactness of $\mathcal{B}$.

\subsection{Band gaps as rule-type incompatibility}

For two adjacent bands $n$ and $n+1$, the maximum of $E_n(k)$ and the minimum of $E_{n+1}(k)$ define the band-gap window:
\[
\Delta E_{n,n+1} = \min_{k \in \mathcal{B}} E_{n+1}(k) - \max_{k \in \mathcal{B}} E_n(k).
\]
If $\Delta E_{n, n+1} > 0$, there is a \emph{band gap}: no eigen-energy of $H$ falls in $(\max E_n, \min E_{n+1})$.

\textbf{Substrate-level reading.} A band gap is the substrate-level statement that \emph{no rule-type-compatible eigenchannel exists at energies in the gap}. The chain at a gap-energy cannot maintain a coherent identity across the periodic substrate: the rule-type structure on adjacent unit cells is incompatible with the chain's would-be propagation at that energy. The chain does not reflect, scatter, or attenuate within the substrate at gap energies --- it \emph{fails to instantiate} a coherent eigenchannel.

This is the substrate-level mechanism behind photonic bandgaps: the substrate's periodic rule-type structure forbids certain energy ranges by incompatibility, not by destructive interference of pre-existing waves.

\subsection{Origin of band gaps: the tight-binding example}

To make the band-gap argument concrete: consider a periodic potential $V(x)$ that is strongly peaked on a sub-lattice of the unit cell, with weak coupling between adjacent peaks. In the tight-binding limit, $u_{n,k}(x)$ is approximately a linear combination of atomic-orbital-like functions localized at each site. The eigen-energies are
\[
E_n(k) = \epsilon_n - 2t_n \cos(ka) + O(t_n^{2})
\]
for nearest-neighbor hopping amplitude $t_n$. The band-width of the $n$-th band is $4|t_n|$, and band gaps appear between bands corresponding to different atomic orbitals when the energy splitting between orbitals exceeds the band-widths. This is a representative case; the structural conclusion (gaps from incompatibility across cells) is general.

\section{Reciprocal Space and Brillouin Zone Geometry}

\subsection{Reciprocal lattice}

The set of translations preserving the substrate's rule-type structure is generated by $T_a$, with periodicity $a$ in real space. The dual lattice in $k$-space is generated by
\[
G \equiv 2\pi/a,
\]
the \emph{reciprocal lattice vector} in one dimension. Two quasi-momenta differing by an integer multiple of $G$ are substrate-levelly identical:
\[
k \sim k + nG \quad \text{for } n \in \mathbb{Z}.
\]
The Brillouin zone $\mathcal{B} = (-G/2, G/2] = (-\pi/a, \pi/a]$ is the fundamental domain --- the set of inequivalent $k$ values.

\subsection{Periodicity of $E_n(k)$}

The band dispersion $E_n(k)$ is periodic in $k$ with period $G$:
\[
E_n(k + G) = E_n(k).
\]
This follows from the equivalence $k \sim k + G$ in the Brillouin-zone identification (\S4.4). Equivalently, $H_{k+G}$ and $H_k$ are unitarily equivalent: define $U_G \psi(x) = e^{-iGx}\psi(x)$, and verify that $U_G H_{k+G} U_G^{-1} = H_k$ on $a$-periodic functions (because $e^{-iGx}$ is itself $a$-periodic when $G = 2\pi/a$).

\subsection{Higher dimensions}

For $d$-dimensional substrates with periodicity in each direction (lattice vectors $\mathbf{a}_1, \ldots, \mathbf{a}_d$), the reciprocal lattice is generated by $\mathbf{b}_i$ satisfying $\mathbf{b}_i \cdot \mathbf{a}_j = 2\pi\delta_{ij}$, and the Brillouin zone is the Wigner--Seitz cell of the reciprocal lattice. The Bloch theorem generalizes straightforwardly:
\[
\psi_\mathbf{k}(\mathbf{x}) = e^{i \mathbf{k}\cdot\mathbf{x}} u_\mathbf{k}(\mathbf{x}), \qquad u_\mathbf{k}(\mathbf{x} + \mathbf{a}_j) = u_\mathbf{k}(\mathbf{x}) \text{ for all } j.
\]
The structural argument is the same: the abelian translation group's irreducible representations are labeled by $\mathbf{k} \in \mathcal{B}$, and the simultaneous eigenchannels of $H$ and the translation operators take the Bloch form.

For the substrate-level account of \emph{photonic Chern insulators} (\S7 and beyond), the relevant case is $d = 2$ (e.g., the synthetic Brillouin zone of the Ch\'enier et al. experiment is two-dimensional). The Brillouin zone is then a torus $\mathbb{T}^2$ --- a closed two-dimensional surface --- which is the setting for the integer-quantization of integrated Berry curvature (the Chern number).

\section{Berry Curvature Across the Brillouin Zone}

\subsection{Setup: the band manifold}

Fix a band index $n$. As $k$ varies over the Brillouin zone $\mathcal{B}$, the periodic factor $u_{n,k}(x)$ varies smoothly in the parameter $k$. The Brillouin zone is the parameter manifold; the Hilbert space of $a$-periodic functions $\mathcal{H}_{\mathrm{cell}}$ is the fibre. The smooth assignment $k \mapsto u_{n,k} \in \mathcal{H}_{\mathrm{cell}}$ defines a section of the parameter-space bundle $E \to \mathcal{B}$ (P-BL-5).

\subsection{The Berry connection}

For each band $n$, define the Berry connection on the Brillouin zone:
\[
\mathcal{A}_n(k) \equiv i \langle u_{n,k} | \partial_k u_{n,k}\rangle_{\mathrm{cell}}.
\]
In one dimension, $\mathcal{A}_n(k)$ is a real-valued function of $k$ (a one-form on $\mathcal{B}$). In two dimensions, $\mathcal{A}_n(\mathbf{k}) = (\mathcal{A}_n^x(\mathbf{k}), \mathcal{A}_n^y(\mathbf{k}))$ is a real-valued vector field on $\mathcal{B}$.

\textbf{Reality of $\mathcal{A}_n$.} Differentiate the orthonormality $\langle u_{n,k}|u_{n,k}\rangle_{\mathrm{cell}} = 1$ with respect to $k$:
\[
\langle \partial_k u_{n,k} | u_{n,k}\rangle_{\mathrm{cell}} + \langle u_{n,k} | \partial_k u_{n,k}\rangle_{\mathrm{cell}} = 0,
\]
so
\[
\langle u_{n,k} | \partial_k u_{n,k}\rangle_{\mathrm{cell}} = -\langle \partial_k u_{n,k} | u_{n,k}\rangle_{\mathrm{cell}} = -\langle u_{n,k} | \partial_k u_{n,k}\rangle_{\mathrm{cell}}^*.
\]
Therefore $\langle u_{n,k}|\partial_k u_{n,k}\rangle$ is purely imaginary, and $\mathcal{A}_n(k) = i\langle u_{n,k}|\partial_k u_{n,k}\rangle$ is purely real.

\subsection{The Berry curvature}

In two dimensions, define the Berry curvature as the curl of the Berry connection:
\[
\Omega_n(\mathbf{k}) \equiv \partial_{k_x} \mathcal{A}_n^y - \partial_{k_y} \mathcal{A}_n^x.
\]
In one dimension, the curvature degenerates to a derivative:
\[
\Omega_n(k) \equiv \partial_k \mathcal{A}_n(k),
\]
though strictly speaking this is the derivative of a one-form rather than a two-form curvature. In one dimension the integrated Berry phase around the circle $\mathcal{B}$ is the only invariant; in two and higher dimensions the curvature is genuinely two-form-valued.

\subsection{Gauge transformations}

The periodic factor $u_{n,k}(x)$ is determined by the eigenvalue equation only up to an $x$-independent phase: $u_{n,k}(x)$ and
\[
\tilde u_{n,k}(x) \equiv e^{i\chi(k)} u_{n,k}(x)
\]
both satisfy $H_k u = E_n(k)\, u$ with the same eigen-energy. Compute the Berry connection in the new gauge:
\begin{align*}
\tilde{\mathcal{A}}_n(k) &= i \langle \tilde u_{n,k} | \partial_k \tilde u_{n,k}\rangle \\
&= i \langle u_{n,k} | e^{-i\chi} \partial_k [e^{i\chi} u_{n,k}]\rangle \\
&= i \langle u_{n,k} | e^{-i\chi} [i (\partial_k \chi) e^{i\chi} u_{n,k} + e^{i\chi} \partial_k u_{n,k}]\rangle \\
&= -\langle u_{n,k}|u_{n,k}\rangle \partial_k \chi + i \langle u_{n,k} | \partial_k u_{n,k}\rangle \\
&= \mathcal{A}_n(k) - \partial_k \chi(k).
\end{align*}
The Berry connection transforms as a $U(1)$ connection one-form: $\mathcal{A}_n \to \mathcal{A}_n - \partial_k\chi$.

\subsection{Gauge-invariance of the curvature}

In two dimensions, the Berry curvature is the curl of the connection. Under a gauge transformation:
\[
\tilde\Omega_n = \partial_{k_x}(\mathcal{A}_n^y - \partial_{k_y}\chi) - \partial_{k_y}(\mathcal{A}_n^x - \partial_{k_x}\chi)
= (\partial_{k_x} \mathcal{A}_n^y - \partial_{k_y} \mathcal{A}_n^x) - 0 = \Omega_n,
\]
using the equality of mixed partials for smooth $\chi$. The Berry curvature is \emph{gauge-invariant unconditionally} (not merely modulo $2\pi$).

\subsection{The Berry phase across the Brillouin zone}

For a closed loop $C$ in $\mathcal{B}$, the Berry phase of band $n$ is
\[
\gamma_n[C] \equiv \oint_C \mathcal{A}_n(\mathbf{k}) \cdot d\mathbf{k}.
\]
In one dimension, $\mathcal{B}$ itself is a circle, and the natural closed loop is the entire BZ. The Berry phase
\[
\gamma_n[\mathcal{B}] = \oint_\mathcal{B} \mathcal{A}_n(k)\, dk
\]
is called the Zak phase. It is gauge-invariant modulo $2\pi$ (as derived for the parameter-space Berry phase: a single-valued gauge transformation $\chi(k)$ cannot change $\gamma_n[\mathcal{B}]$, because $\chi(k+G) - \chi(k)$ must be an integer multiple of $2\pi$ for $e^{i\chi}$ to be single-valued on the circle).

\subsection{Stokes' theorem on the Brillouin zone (2D case)}

In two dimensions, for a simply-connected region $S \subset \mathcal{B}$ with boundary $\partial S = C$:
\[
\oint_C \mathcal{A}_n \cdot d\mathbf{k} = \iint_S \Omega_n\, d^{2}k,
\]
by the standard Stokes' theorem on $\mathbb{R}^2$ (which extends to the torus in regions where the gauge can be chosen smoothly).

\textbf{Closed surface integral.} For a \emph{closed} surface $\Sigma$ (e.g., the entire 2D Brillouin zone, which is topologically a torus), the integrated curvature is
\[
C_n \equiv \frac{1}{2\pi} \iint_\mathcal{B} \Omega_n\, d^{2}k,
\]
and $C_n$ is the \emph{Chern number} of the $n$-th band.

\subsection{Integer quantization of the Chern number --- flagged OPEN}

\textbf{Statement.} When $\mathcal{B}$ is a closed manifold (a torus in 2D, a sphere in 3D parameter spaces such as the Bloch sphere), the integrated Berry curvature
\[
\iint_\mathcal{B} \Omega_n\, d^{2}k = 2\pi C_n, \qquad C_n \in \mathbb{Z}.
\]

\textbf{Proof sketch.} The integer-quantization is forced by the fact that the parameter-space bundle is a $U(1)$ principal bundle over a closed manifold. The transition functions between local trivializations on overlapping patches must be $U(1)$-valued and single-valued on the overlaps. The integrated curvature counts the winding number of these transition functions around the boundaries of the patches --- which is an integer by single-valuedness.

\textbf{Substrate-level reading.} The integer-quantization is the statement that the rule-type bundle's holonomy around a closed surface is a topological invariant: the global participation-rule twist, when wrapped around a closed manifold, can only take integer multiples of $2\pi$.

\textbf{Status: OPEN.} A full substrate-level proof of integer-quantization is not given here. The proof inherits from standard differential-topology arguments (the Chern theorem on principal $U(1)$ bundles); a substrate-level account that would make integer-quantization FORCED rather than INHERITED is a candidate for a follow-on Chern-quantization walkthrough.

For the present walkthrough's scope, integer-quantization is FORM-FORCED-INHERITED at the level of the standard differential-topology proof, with a substrate-level reading layered on top.

\section{Substrate-Level Reading}

The objects derived in \S\S3--7 admit the following substrate-level dictionary:

\begin{center}
\small
\begin{tabular}{@{}p{5cm}p{9cm}@{}}
\toprule
Standard band-theory object & Substrate-level meaning \\
\midrule
Periodic potential $V(x+a) = V(x)$ & Periodic rule-type structure (P-BL-1) \\
Translation operator $T_a$ & Substrate-level rule-shift by one period (P-BL-4) \\
Commutation $[H, T_a] = 0$ & Chain evolution respects substrate's discrete translation symmetry \\
Bloch eigenchannel $\psi_k(x) = e^{ikx}u_k(x)$ & Eigenchannel of the joint $H{+}T_a$ algebra \\
Quasi-momentum $k$ mod $2\pi/a$ & Parameter-space label for the substrate's translation-symmetry-irreducible representation \\
Brillouin zone $\mathcal{B}$ & Fundamental parameter-space domain; circle (1D) or torus (2D) \\
Band $E_n(k)$ & Discrete spectrum of substrate-level eigen-energies indexed by quasi-momentum \\
Band gap & Substrate-level rule-type incompatibility; chain fails to instantiate coherent identity \\
Berry connection $\mathcal{A}_n(k)$ & Rule-type connection on the parameter-space bundle over $\mathcal{B}$ \\
Berry curvature $\Omega_n(k)$ & Curvature of the rule-type connection in parameter space \\
Chern number $C_n$ & Integer-valued global topological invariant of the band's rule-type bundle over $\mathcal{B}$ \\
\bottomrule
\end{tabular}
\end{center}

\subsection{Why band gaps are absolute, not soft}

A standard wave-mechanical reading of band gaps invokes destructive Bragg interference: forward-propagating and backward-propagating waves cancel at the gap-band edge. This reading does not explain why the gap is absolute --- why no eigenchannel exists in the gap, rather than merely a strongly-attenuated one.

The substrate-level reading: at gap energies, the chain's would-be Bloch eigenchannel ($\psi_k(x) = e^{ikx}u_k(x)$ with the gap energy as eigenvalue) requires a periodic factor $u_k$ that does not exist as a square-integrable solution of the $H_k$ eigenvalue problem. The chain cannot align its identity with the substrate's periodic rule-type structure at these energies. The gap is structural --- a substrate-level rule-type incompatibility, not a dynamical interference effect. (Bragg interference is the wave-mechanical \emph{coarse-grained image} of the substrate-level rule-type incompatibility.)

\subsection{Why the Brillouin zone is a closed manifold}

The substrate-level statement that translations by $a$ leave the rule-type structure invariant (P-BL-1) forces the quasi-momentum $k$ to be defined modulo $G = 2\pi/a$. This identification makes the BZ topologically a circle in 1D, a torus in 2D. The closedness of the BZ as a manifold is what makes the integrated Berry curvature (Chern number) a meaningful integer-valued topological invariant --- there are closed loops and closed surfaces to integrate over.

\subsection{Why the band index is discrete}

The discreteness of the band index $n$ is forced by the spectrum of $H_k$ on the unit cell being discrete (\S5.3). The unit cell is a compact domain with periodic boundary conditions; the kinetic-plus-potential operator on a compact domain has a discrete spectrum by standard spectral theory. The substrate-level reading: each band is a distinct rule-type identity that the chain can commit to, indexed by an integer.

\section{What's Forced, What's Inherited, What's Open}

\subsection*{What is FORCED by the substrate ontology}

\begin{itemize}[noitemsep]
  \item \textbf{Bloch's theorem} $\psi_k(x) = e^{ikx}u_k(x)$ with $u_k(x+a) = u_k(x)$ is FORCED by the commutation $[H, T_a] = 0$ + unitarity of $T_a$ + simultaneous-eigenbasis structure for commuting Hermitian operators (\S4).
  \item \textbf{Brillouin zone identification} $k \sim k + 2\pi/a$ is FORCED by the discrete translation symmetry of the periodic substrate (\S4.4).
  \item \textbf{Band eigenvalue problem} $H_k u_{n,k} = E_n(k)u_{n,k}$ on the unit cell is FORCED by reduction of the full Schr\"odinger problem via the Bloch ansatz (\S5.1).
  \item \textbf{Discrete band structure} is FORCED by the compactness of the unit cell + discreteness of the spectrum of Hermitian operators on compact domains (\S5.3).
  \item \textbf{Band gaps as rule-type incompatibility} are FORCED whenever the unit-cell eigenvalue problem $H_k u = E u$ has no $a$-periodic square-integrable solution at a given $E$ for any $k \in \mathcal{B}$ (\S5.4).
  \item \textbf{Berry connection} $\mathcal{A}_n(k) = i\langle u_{n,k}|\partial_k u_{n,k}\rangle$ on the BZ is FORCED by smooth $k$-dependence of the band + orthonormality of periodic factors (\S7.2).
  \item \textbf{Gauge-invariance of Berry curvature} $\Omega_n$ is FORCED by the differential-form structure of the connection ($d^2 = 0$ for smooth $\chi$) (\S7.5).
  \item \textbf{Periodicity of $E_n(k)$} with reciprocal-lattice period $G$ is FORCED by the BZ identification + unitary equivalence of $H_{k+G}$ and $H_k$ (\S6.2).
\end{itemize}

\subsection*{What is FORM-FORCED-INHERITED (and re-derived inside this document)}

\begin{itemize}[noitemsep]
  \item \textbf{Spectral theorem for commuting Hermitian operators} (\S3.3): re-derived briefly via simultaneous diagonalization.
  \item \textbf{Unitarity of translation operator} $T_a^\dagger T_a = I$ (\S P-BL-4): re-derived from change of variable in the inner product.
  \item \textbf{Differentiation-translation commutation} $T_a\partial = \partial T_a$ (\S3.2): re-derived from the chain rule.
  \item \textbf{Compactness $\Rightarrow$ discrete spectrum} for Hermitian operators on the unit cell (\S5.3): standard spectral theory; cited rather than re-proven.
  \item \textbf{Stokes' theorem} in 1D and 2D (\S7.7): standard; cited.
  \item \textbf{Chain rule, integration by parts, change of variables}: standard calculus; used inline.
\end{itemize}

\subsection*{What remains OPEN}

\begin{itemize}[noitemsep]
  \item \textbf{Integer-quantization of the Chern number} (\S7.8). Inherited at the level of the standard differential-topology proof. A substrate-level account that makes integer-quantization FORCED (rather than INHERITED) is a candidate for a follow-on Chern-quantization walkthrough.
  \item \textbf{Tight-binding limit derivation.} The \S5.5 example invoked the tight-binding form without derivation. The full substrate-level account of the tight-binding limit (Wannier functions, exponential localization) is OPEN; FORM-FORCED-INHERITED at the standard solid-state level.
  \item \textbf{Beyond simple periodic substrates.} Quasi-periodic, disordered, and incommensurate substrates require generalization. OPEN.
  \item \textbf{Time-reversal-broken substrates and the topology shift.} The Bloch theorem itself does not require time-reversal symmetry; band-topology classification depends on which symmetries are preserved. The Haldane-model construction (which breaks time-reversal symmetry) requires a separate analysis.
  \item \textbf{Non-Abelian Berry connection} for degenerate bands (band crossings, Wilczek--Zee structure). OPEN; pair-walkthrough candidate.
\end{itemize}

\section{What This Argument Establishes}

This walkthrough establishes the following exact claims:

\textbf{Claim 1.} Bloch's theorem $\psi_k(x) = e^{ikx}u_k(x)$ with $u_k(x+a) = u_k(x)$ is FORCED by the substrate ontology when a chain propagates through a periodic rule-type substrate. The mechanism is the commutation $[H, T_a] = 0$ + unitarity of $T_a$ + simultaneous-eigenbasis structure.

\textbf{Claim 2.} The Brillouin zone $\mathcal{B} = (-\pi/a, \pi/a]$ with endpoints identified is FORCED as the fundamental parameter-space domain of the substrate's discrete translation symmetry. Topologically, $\mathcal{B}$ is a circle in 1D, a torus in 2D.

\textbf{Claim 3.} The band eigenvalue problem $H_k u_{n,k} = E_n(k)u_{n,k}$ on the unit cell is FORCED by reduction via the Bloch ansatz. The discrete band structure $\{E_n(k) : n \in \mathbb{N}\}$ at each $k$ is FORCED by the compactness of the unit cell and the spectral theorem.

\textbf{Claim 4.} Band gaps are FORCED at energies for which no $a$-periodic square-integrable eigenchannel of $H_k$ exists for any $k \in \mathcal{B}$. The substrate-level statement: a band gap is rule-type incompatibility --- the chain cannot align its identity with the periodic rule structure at gap energies. Bragg interference is the coarse-grained wave-mechanical image of this incompatibility.

\textbf{Claim 5.} The Berry connection $\mathcal{A}_n(k) = i\langle u_{n,k}|\partial_k u_{n,k}\rangle$ is FORCED on the Brillouin zone by smooth $k$-dependence + orthonormality. The Berry curvature $\Omega_n(\mathbf{k}) = \partial_{k_x}\mathcal{A}_n^y - \partial_{k_y}\mathcal{A}_n^x$ (in 2D) is gauge-invariant.

\textbf{Claim 6.} The integrated Berry curvature over a closed Brillouin zone, $C_n = (1/2\pi)\iint_\mathcal{B}\Omega_n\, d^2k$, is integer-valued. This is FORM-FORCED-INHERITED from standard differential topology; substrate-level account flagged OPEN.

\textbf{Claim 7 (negative).} No new substrate primitives are required. Bloch's theorem, the band structure, and the Berry connection on the BZ all follow from composition of: periodic participation rule (P-BL-1), identity alignment under spatial translation (P-BL-2), channel propagation through periodic gradients (P-BL-3), translation operator (P-BL-4), and parameter-space bundle over the BZ (P-BL-5), plus standard linear-algebra and calculus machinery re-derived as needed.

\textbf{Claim 8 (scope-limit).} This walkthrough does not derive: integer-quantization of the Chern number from substrate primitives (inherited from differential topology); tight-binding-limit Wannier-function structure; beyond-simple-periodic substrates; full topological classification under symmetry; or non-Abelian Berry connection for degenerate bands.

\textbf{The unified statement.} A periodic rule-type substrate forces the chain's eigenchannels into Bloch form, and the chain's parameter-dependent identity across the resulting Brillouin zone carries a rule-type connection (Berry connection) and curvature (Berry curvature) that encode the band's geometric structure. Band gaps are substrate-level rule-type incompatibilities; the Chern number is the global topological invariant of the band's rule-type bundle over the closed Brillouin zone.

\section*{References}

\begin{itemize}[leftmargin=*,itemsep=2pt]
  \item Bloch, F. \emph{\"Uber die Quantenmechanik der Elektronen in Kristallgittern.} Z. Phys. \textbf{52}, 555 (1929).
  \item Ashcroft, N. W., Mermin, N. D. \emph{Solid State Physics.} Harcourt (1976).
  \item Kittel, C. \emph{Introduction to Solid State Physics.} Wiley (8th ed., 2004).
  \item Simon, B. \emph{Holonomy, the quantum adiabatic theorem, and Berry's phase.} Phys. Rev. Lett. \textbf{51}, 2167 (1983).
  \item Thouless, D. J., Kohmoto, M., Nightingale, M. P., den Nijs, M. \emph{Quantized Hall conductance in a two-dimensional periodic potential.} Phys. Rev. Lett. \textbf{49}, 405 (1982).
  \item Xiao, D., Chang, M.-C., Niu, Q. \emph{Berry phase effects on electronic properties.} Rev. Mod. Phys. \textbf{82}, 1959 (2010).
  \item Nakahara, M. \emph{Geometry, Topology, and Physics.} IOP Publishing (2003).
\end{itemize}

\end{document}
